\exercisesection{3}

\begin{enumerate}
\def\labelenumi{\arabic{enumi}.}
\item
  Under the hypotheses of Proposition 1 of no. 1, show that every prime ideal \(\mathfrak{p} \in \operatorname{Ass}_f(\mathbf{M})\) (§ 2, Exercise 17) is graded and is a minimal element of the set of prime ideals containing the annihilator of a homogeneous element of M. (Show first that the largest graded ideal \(\mathfrak{p}'\) contained in a prime ideal \(\mathfrak{p}\) of A is prime, observing that, if the product of two elements \(x, y\) belongs to \(\mathfrak{p}'\) , so does the product of their homogeneous components \(\neq 0\) of highest degree. Note then that, if a prime ideal contains the annihilator of an element \(z \in \mathbf{M}\) , it also contains the annihilator of at least one of the homogeneous components \(\neq 0\) of \(z\) .)
\item
  In a graded ring A of type \(\Delta\) (where \(\Delta\) is torsion-free) let p be a minimal element of the set of prime ideals of A, which is necessarily graded (Exercise 1). Show that for every homogeneous element \(t \in p\) there exists a homogeneous element \(s \notin p\) and an integer n \textgreater{} 0 such that \(st^{n} = 0\) . Deduce generalizations of Propositions 3, 4 and 5 to graded rings which are not necessarily Noetherian (with the definitions of §2, Exercises 12 and 20).
\item
  Let A be a graded ring of type \(\Delta\) (where \(\Delta\) is torsion-free) and M a strongly Laskerian finitely generated graded A-module ( \(\S 2\) , Exercise 28). Show that the submodules \(Q(\mathfrak{p}) = Q(\mathfrak{p}, e_{\mathfrak{p}}(\mathbf{M}))\) (where p runs through \(\operatorname{Ass}_{f}(\mathbf{M})\) ) defined in Exercise 4 of \(\S 2\) are graded submodules.
\end{enumerate}

